\section{PopQA Extension: The Recurring-Memory Result Generalizes}

Table~\ref{tab:j-popqa} shows the broader journal PopQA run. The central observation is that the recurring-memory story survives broader evaluation and remains stable under the calibrated public-track fix. On \texttt{flan-t5-base}, both the full and calibrated routers reach `0.290` quality against `0.290` for retrieval-only while reducing retrieval calls from `100.0` to `71.0`. On \texttt{flan-t5-small}, both routers reach `0.130` versus `0.120` for retrieval-only while still replacing a meaningful fraction of retrieval with memory use.

\begin{table}[h]
\centering
\small
\begin{tabular}{llrrrrr}
\toprule
Model & Mode & Quality & 95\% CI & Retrieve & Read & Write \\
\midrule
flan-t5-base & always\_retrieve & 0.290 & [0.290, 0.290] & 100.0 & 0.0 & 0.0 \\
flan-t5-base & router & 0.290 & [0.290, 0.290] & 71.0 & 73.0 & 48.0 \\
flan-t5-base & router\_calibrated & 0.290 & [0.290, 0.290] & 71.0 & 73.0 & 48.0 \\
flan-t5-base & router\_no\_v\_future & 0.130 & [0.130, 0.130] & 42.0 & 22.0 & 20.0 \\
flan-t5-small & always\_retrieve & 0.120 & [0.120, 0.120] & 100.0 & 0.0 & 0.0 \\
flan-t5-small & router & 0.130 & [0.130, 0.130] & 72.0 & 75.0 & 50.0 \\
flan-t5-small & router\_calibrated & 0.130 & [0.130, 0.130] & 72.0 & 75.0 & 50.0 \\
flan-t5-small & router\_no\_v\_future & 0.090 & [0.090, 0.090] & 26.0 & 6.0 & 5.0 \\
\bottomrule
\end{tabular}

\caption{Broader seeded PopQA result from \texttt{artifacts/journal\_popqa\_broad\_100\_calibrated\_v1/aggregate\_table.md}.}
\label{tab:j-popqa}
\end{table}

This broader run is consistent with the frozen recurring-heavy slice rather than contradicting it. The base-model result is especially important because it shows that the router remains close to the retrieval baseline while materially changing the operation mix. The ablation dependence also survives: removing reads, writes, or delayed utility sharply degrades performance. In other words, the PopQA result is not a lucky single-slice artifact anymore; it is a repeated recurring-memory pattern.
